\documentclass[eprint,
approc,		%heading for Acta Polytechnica CTU Proceedings
]{actapoly}

\usepackage[utf8]{inputenc}
%--------------my packages ------------>
\usepackage{float}
\usepackage{soul}
\graphicspath{{figures/}} % path to folder with figures
\usepackage{subcaption}
\usepackage{rotating}
\usepackage{dcolumn}
\usepackage{tablefootnote}
\usepackage{listings}
\usepackage{xcolor}
%\usepackage{gensymb} % for \textdegree
\usepackage{tabularx}
\usepackage{float}
\usepackage{chngcntr}
\counterwithin*{equation}{section}
\usepackage[super]{nth}
\usepackage{longtable}
\usepackage{tabularx}
\usepackage{multirow}
\usepackage{adjustbox}
\graphicspath{{figures/}} % path to folder with figures
\usepackage{graphicx}
%--------------my packages ------------<

\newcommand*\printft[1]{\textsc{\lowercase{#1}}}

\begin{document}

\title[Classification Techniques using Machine Learning Algorithms...]
{Classification Techniques using Machine Learning Algorithms for Information Extraction from Landsat-8-9 OLI/TIRS Sensor Data}

\correspondingauthor[P. Lemenkova]{Polina Lemenkova}{my}{polina.lemenkova2@unibo.it}

\institution{my}{University of Bologna, Department of Biological, Geological and Environmental Sciences, Via Irnerio 42, 40126 Bologna, Emilia-Romagna, Italy}

\begin{abstract}
Methods for image processing are essential for geo-information extraction. This work presents the use of Machine Learning (ML) for image processing. The aim is to recognise landscape patches using ML classification in GRASS GIS modules. The algorithms are integrated from Python library Scikit-Learn. The data included a time series of Landsat scenes. The workflow of image processing includes image classification for automatic detection of land categories. The algorithms SupportVectorMachine, MLPC of ANN, Decision Tree Classifier, and RandomForest were employed to identify land patterns by computer vision. The topology of  patches was detected by ML that considers the differences among spectral reflectance of pixels. ML recognises texture and spectral features to measure the similarity of neighbouring patches while generating random decision trees. This study shown the advantages of ML for remote sensing and GIS.
\end{abstract}

\keywords{artificial neural network, decision tree, random forest, support vector machine, spatial data} 

\maketitle

\section{Introduction}

Recently, the use of machine learning (ML) has been in the centre of attention recently for the challenge of Earth observation data processing \citep{9395808,10041606,Lemenkova202320,8547145}. ML has become available approach in cartography along with the advances in programming and these methods are increasingly used in various tasks of environmental monitoring and mapping \citep{10142835,jimaging9050098,9850602}. \hl{Their use }ensures automation of satellite data processing, supports operative monitoring using time series and has an important effect on accuracy and effectiveness of data processing for cartographic visualization \citep{jmse12081279}. Thus, high level of automation in data processing achieved by ML approaches makes it possible to map land cover types rapidly using computer vision approaches applied to satellite images.

\hl{Early detection of landscapes dynamics needs a timely recognition of subtle changes in the forests and involves Earth monitoring. One of the key approaches for identification landscape dynamics in land cover types is the analysis of changes and the rate of exposure of landscapes to climate and environmental effects using multi-temporal geospatial data. The effectiveness of the remote sensing (RS) data for environmental forest monitoring is explained by the optical properties of satellite images: the features of land cover types on the Earth have different spectral reflectance values (Dinh Duong, 2018; Lemenkova, 2023; Gao et al., 2006). This enables to identify them using digital signatures of the satellite images. Time series analysis of such scenes enables to detect changes and recognise modified patches within the landscape structure (Lemenkova \& Debeir, 2023; Babadi Ataabadi et al., 2025; Wu et al., 2025). Extraction of these features from the RS signals recorded during the Earth monitoring from Landsat space missions provides an essential tool to detect landscape dynamics. }

The state-of-the-art methods of image classification in traditional GIS methods have certain limitations which can be summarised as follows. Firstly, since image classification is generally based on spectral reflectance of pixels, some cases may take place. This is cased by the ground properties of pixels which are not representing real land surface objects and strongly depend on the resolution of satellite images \citep{6351712,Lemenkova202322}. As a result, objects which have very similar spectral properties can be erroneously assigned to the same land cover class. Secondly, object topology is often avoided in when interpreting the imagery. To avoid this drawback, some studies analyse spatial relationships between the objects using Object Based Image Analysis (OBIA) techniques \citep{LUCIANO2019127,PHIRI2018170,DEWALI20233109,MAO202011}. 

Nevertheless, this is a different approach that requires specific software and tools. Thirdly, the inherited properties of the land surface such as texture, curvature of patches, context of disposition of several objects representing the scene, shape, form and structure of objects is difficult to properly identified using the traditional image classification techniques in standard GIS \citep{OUSMANOU2024100239,ZHANG2023129590,BAKR2010592}. Moreover, increased variability of images with different properties and resolution require more advanced methods compared to the traditional approaches of image analysis that use pixel-based classification algorithms. This lead to the issues in accuracy and precision of image classification process \citep{Fisher,McKee}.

In view of these limitations, the use of the novel ML methods for RS data processing and programming algorithms for classification approaches has recently drawn much attention \citep{Ayhan}. Among them, Python algorithms present a powerful instrument for image processing that have been investigated in a number of works \citep{9884244,9883588,Lemenkova202227,9006205,7724471}. The RF algorithm embedded in the mapping of crop fields using Landsat imagery was thoroughly described by \citep{LUCIANO2021106063}. \citep{AHMED201589} proposed an integrated application of the RF algorithm based on the multi-source data combining LiDAR and Landsat where connected components in the RS data are particularly well suited for complex forest canopy analysis using data on cover and height. Diverse examples of the RF algorithm applications in various cases of RS data processing include crop mapping for food and agricultural research \citep{OLIPHANT2019110,TELUGUNTLA2018325}, hazard risk assessment using fire mapping \citep{COLLINS2020111839}, urban sprawl analysis \citep{GHOSH201431} and environmental monitoring of land cover types \citep{MELVILLE201846,GRINAND201368}. 

More research projects have demonstrated the potential of the RF approach of ML to image processing \citep{COLLINS2020111839,WEI2020112005,WANG202294}. Nevertheless, it still suffers from some limitations regarding its application in the GRASS GIS. For example, RF algorithm needs to be better exploited using GRASS GIS scripting modules in view of the current lack of a practical framework and hands-on applications in case studies. There have been limited case studies on the use of GRASS GIS using scripts for image processing in environmental applications \citep{STRIGARO201668,geomatics5010005,ROCCHINI2017166,7326254}. 

In further examples, \citep{7326254} evaluated the use of GRASS GIS modules for marine geological mapping using MODIS/Aqua RS data in coastal areas. The approach presented by \citep{6407912} reports the links between GRASS GIS and Python and performs relatively well for environment monitoring using satellite imagery but the proper scripts of GRASS GIS are not reported since the major point was using Python for launching GRASS GIS. Moreover, few experimental studies have thoroughly evaluated the GRASS GIS modules for diverse ecological applications and mapping \citep{lemenkova2024exploitation}. However, the method of using RF by the GRASS GIS is not reported in the existing available works and needs evaluation.

\section{Materials and Methods}

\hl{This paper evaluates image classification methods using training algorithm of the Machine Learning (ML) techniques of GRASS GIS including the embedded Scikit-Learn ML library of Python (Lemenkova, 2025a). The ML approach with different classifiers present the example of advanced cartographic methods of image classification. Land cover types around central Apennines, Italy, were identified using classification of multispectral Landsat 8-9 OLI/TIRS images collected in different seasons using modules of GRASS GIS 'r.learn.predict', 'r.random', 'r.learn.train', 'r.category', 'i.group'. This enabled to monitor landscape dynamics during the estimated period. Comparative analysis revealed variations in the surface of basin from 2018 to 2023 and evaluate land cover dynamics. The series of maps made using GRASS GIS scripts demonstrated the opportunities of the ensemble learning tools to image processing. Scripting approach enables accurate image classification for environmental analysis of landscape dynamics and supports cartographic visualization. The ML of image classification proposes accurate tool for monitoring landscapes changes.}

To classify RS data, two types of approaches were carried out. The first is unsupervised classification based on the MaxLike approach using 'i.cluster' and 'i.maxlik' and clustering which extracted Digital Numbers (DN) of cell feature based on spectral reflectance of signals. The second is supervised classification using several ML methods, technically realised using GRASS GIS scripts.

\subsection{Study area: A brief overview of the environmental settings}

\hl{The topography of the Central Apennines is characterised by diversity of accumulated and depleted landforms that resulted from the millennial history of geologic and soil formations of Apennines, Figure} \ref{fig:soil}. \hl{The area of central Apennines is notable for high forest cover and high biodiversity according to the soil types. The dominating vegetation species include Silver fir, Norway spruce, Beech, Black pine, Birch). Examples of deciduous forests include such types as mixed oak \emph{Quercus [L.]}, hornbeam \emph{Carpinus [L.]}, and beech forests \emph{Fagus sylvatica [L.]}, in addition to vegetation typical for woodlands situated near streams, e.g. willow \emph{Salix [L.]} and aspen \emph{Populus tremula [L.]} (Ballelli, \& Allegrezza, 2016). }

The distribution of dominating classes, such as European beech \emph{Fagus sylvatica [L.]} forests and other land cover classes is controlled by a variety of factors including the topographic structure, environmental and climate settings. The variability of environmental factors affects habitat structure and influences species composition of forests and vascular plants. In valleys, such areas are expanding and provide a valuable resource for locals, e.g., the production of fruit and lumber.  The variety of landforms and topographic structure created conditions for diverse soil types. This provides essential ecosystem services to habitats and includes, for instance, such precious forests as mixed silver fir \emph{(Abies alba Mill [L.])}. 

Conifers, including silver fir \emph{Abies alba [L.]}, common spruce \emph{(Picea abies [L.])}, Austrian pine \emph{(Pinus nigra [L.])}, and cedar species \emph{(Cedrus elegans [L.])}, are frequent in plantations that were developed starting in the middle of the 20th century. In the past, areas of mountain forests were used as pastures, for agriculture fields, or also exploited as fuelwood and timber. Among deciduous species dominates European beech \emph{(Fagus sylvatica [L.])}, followed by sweet chestnut \emph{(Castanea sativa [L.])}, whose historical range was anthropogenically expanded for providing valuable resources for locals. These types of forest and plants are adjusted to the specific soil typical for central Italy. The intensive use of these resources persisted until the middle of the 20th century despite certain variations in the use of mountain forests caused by natural fluctuations in the population growth. 

\begin{figure}[H]
\centering
	\includegraphics[width=7cm]{Soil.jpg}
\caption{Soil types in Italy according to the Food and Agricultural Organization (FAO) dataset: 10 major categories generalised from the database with thematic maps showing soil regions of Italy, scaled 1:5,000,000. Software: QGIS.}
\label{fig:soil}
\end{figure}

%---------------------------->
\subsection{Conceptual approach}
The GIS that have included ML modules to select from perform better for image processing compared to simplified programs. Amongst the types of such available tools, Geographic Resources Analysis Support System (GRASS) GIS provides the options of ML for RS data processing that are based on the Python's library Scikit-Learn \citep{scikitlearn}. For instance, Random Forest (RF) learning classifier, linear classifiers that include Support Vector Machine (SVM), logistic regression and GradientBoostingClassifier can be used using GRASS GIS for satellite image processing. 

These options are not only intended to advance the image processing workflow using powerful tools of Python embedded in GRASS GIS thus ensuring the quality and accuracy of image classification but also use a variety of mathematical approaches that can be finely tuned to diverse purposes of image processing. While GRASS GIS also has a wide variety of traditional modules for processing vector and raster types of geospatial data, such as tools for computations of land cover types, geomorphological and hydrological modelling or analytical data visualisations, these options differ in that they predominantly adjusted to cartographic data using internal modules rather than Python algorithms. 

\hl{State-of-the-art image classification approach depends on the regional landscape characteristics and specific property of land cover types. This approach was created to enhance change detection categorization for Italian landscapes and can be further extended to other regions. Landscape heterogeneity, surface ruggedness, slope curvature and regional spatial complexity of land surface may affect the degree of fragmentation. In contrast, the ML algorithms are more efficient and straightforward in cartographic implementation in comparison with traditional ways of image classification. The main goal and scope of this study which limits the approach was finding the best ML algorithm for spatial map aggregation. Here, we tested techniques of data integration and processing using several ML methods for image classification for the region of Central Apennines, Italy.}

\subsection{Gaps and limitations}

\hl{Current project contributes to fill in the existing gaps and limitations related to the environmental monitoring of Italy. First, there is insufficient vegetation mapping of the area, hence, this study contributes through integrating new data from the existing collections of the satellite images. The integration of the new methods, data and literature search with newly collected data supports ecological monitoring of Alps. Second, comparing to the VHR (very high resolution) images with no granted free access, Landsat images contribute to the extraction of the environmental information using open access. Third, permission of the biodiversity survey is often restricted in the natural reserves and protected areas. In such case, remote sensing data enable to perform Earth's observation using non-destructive methods and open access tools.}

\hl{Finally, the mismatch between the project calendar and phenological seasons of vegetation growth often create difficulties in direct field observations for environmentalists and land planners. For example, this includes the impairing the preparation of the fieldwork and rearranging of the survey activities across two consecutive calendar years. In contrast to this, RS data modelling is independent from the seasonal factors and practical challenges of the environmental fieldwork. To fill in these gaps, current study contributes to better understanding of impacts of climate change on vegetation in Central Italy through RS data processing using ML methods. To this end, the key framework is borrowed methodologically from several domains: satellite image processing, machine learning algorithms, and environmental analysis for risk assessment as complex concept of landscapes' exposure and vulnerability to climate effects. }

\hl{The application of the environmental analysis to RS data using ML methods is relatively new concept with few applications due to the complexity of the workflow. Hence, this multi-disciplinary approach might be configured as a new tool for creating environmental predictions of vegetation response to climate change. In the future applications, it can be used or adopted by landscape managers to prioritise the conversation of the plant communities that are vulnerable to external climate effects (such as mountainous ecosystems of Italy) and likely to experience higher-than-average rates of climate change.}

\subsection{Workflow}
\hl{This paper aims to demonstrate the effective use of the ML method of image processing which was pitted against the other traditional classification methods such as k-means Clustering and Maximum Likelihood (MaxLike) of GRASS GIS. The use of such advanced solution helps in providing automated tools of satellite image classification. }The brief description of the workflow is as follows. After data import and preprocessing, the images were first classified using the 'MaxLike' approach of the unsupervised classification which uses clustering (Figure \ref{fig04}). The numerical results of the classification are stored using the 'reportfile' function which contains the statistical parameters of cluster classes obtained from the classified images, as follows: class size, separation, percent convergence, number of iterations, means and standard deviations for Landsat bands and class separability matrix. 

The training for the ML was obtained from the classification and applied as seed dataset. The standard and mean deviation of the partition were calculated and store din the .txt files for each image. Training cluster maps were computed for 10 classes of land cover types over the study area for each scene and used as training data for the ML approach. The outcomes are validated with computed reject threshold results using chi square test to determine levels of confidence of correctly classified pixels on raster scenes\hl{. }Afterwards, the images were classified accordingly using four algorithms of ML\hl{. }Each Landsat scene of the 8 images comprising the dataset was processed and the maps were visualised using these data: images on 2018, 2019, 2022 and 2023 for spring and autumn periods accordingly. 

%---------------------------->
\subsection{Techniques}

Several ML algorithms embedded in the GRASS GIS have been used in this study to classify Landsat 8-9 OLI/TIRS multispectral satellite images based on certain feature combinations. The area is covering the region of central Apennines in a regional scale. Four ML methods of image processing were employed including the following ones: 1) Support Vector Machine (SVM), 2) Multi-layer Perceptron (MLP) Classifier of the Artificial Neural Network (ANN), 3) Decision Tree Classifier (DTC), and 4) Random Forest (RF). These methods are based on different mathematical algorithms embedded in the GRASS GIS. We used these approaches for image classification to detect changes in land cover types over the study area in central Italy and map them during the 5-year period of 2018 to 2023. We tested these algorithms in order to compare them and define the appropriate classifier. Defining optimal classification approach supports the accurate discernibility of the land cover features, and to map the variability of landscapes over time.

\subsection{Dataset}

Major challenges facing a successful classification of the RS data are cloudiness of the satellite images, coverage, data availability and resolution which enable to accurately detect the changing features over the landscapes. Therefore, we selected the quality satellite data with low cloudiness and optimal coverage of the study region. Specifically, we used the moderate resolution multispectral satellite images obtained from the Landsat Collection Level-2. The essential characteristics and key feature attributes of the Landsat 8-9 OLI/TIRS images\hl{.} 

Technically, the data were obtained from the United States Geological Survey (USGS) which presents a collection of the diverse geospatial datasets in the open source repository. Specifically, the sensors used are Landsat 8-9 Operational Land Imager (OLI) and Thermal Infrared Sensor (TIRS). \hl{These sensors were chosen because they }provide the images of Earth's surface and present the collection of data suitable for monitoring landscape dynamics visible on the Earth's surface from space. Such images enable to evaluate trends in land cover changes over time when comparing multi-temporal datasets. The row/path of the Landsat swath is 30/191 for all the scenes, collected in Nadir. The images are projected automatically in the UTM Zone 33 for Italy in WGS84 Ellipsoid/Datum. 

\subsection{Algorithms}

The satellite image processing was performed using GRASS GIS software using the algorithms described below and adopted from the existing works \citep{Lemenkova202217,Lemenkova202313}. The details are given on both theoretical foundations (mathematical base) and practical implementation (essential snippets of scripts of GRASS GIS) of these methods. 

%---------------------------->
\subsubsection{Maximum-likelihood discriminant analysis classifier}

Theoretical base of the algorithm of maximum likelihood (MaxLike) estimation consists in the following. This method pursues the statistical and evaluation goals and is widely used in the remote sensing data analysis. This method aims at evaluation of likelihood that cells within the satellite image belong to the target group. This method\hl{ partitions the} pixels into\hl{ }clusters\hl{. }Technically, it is done using data on DNs of the pixels that vary in values according to different land cover types\hl{, }e.g., water, forest, urban areas, meadows etc, which have distinct spectral reflectance that can be recognised using computer vision algorithms. Similar spectral reflectance within the observed dataset are grouped into the same cluster and vice versa.\hl{ }Practical implementation of this method in GRASS GIS consists in using 'i.cluster' and 'i.maxlik' modules as described in the technical script below. The results of the image processing performed using this method are presented in Figure \ref{fig04}. 

\begin{footnotesize}
\begin{verbatim}
g.region raster=L_2019_01 -p
i.group group=L_2019 subgroup=res_30m input=L_2019_01,\
<...>,L_2019_07
# Clustering: generating signature file and report using \
k-means clustering algorithm
i.cluster group=L_2019 subgroup=res_30m \
signaturefile=cluster_L_2019 
  classes=10 reportfile=rep_clust_L_2019.txt --overwrite
# Classification by i.maxlik module
i.maxlik group=L_2019 subgroup=res_30m \
signaturefile=cluster_L_2019
  output=L_2019_clusters reject=L_2019_cluster_reject
r.colors L_2019_clusters color=roygbiv
# Mapping
g.region raster=L_2019_01 -p
d.mon wx0
d.vect isolines color='100:93:134' width=0
d.rast L_2019_clusters
d.grid -g size=00:30:00 color=grey width=0.1 \
fontsize=16 text_color=grey
d.legend raster=L_2019_clusters title="Clusters 2019" \
title_fontsize=19 font="Helvetica" \
	fontsize=17 bgcolor=white border_color=white
d.out.file output=Italy_2019 format=jpg --overwrite
\end{verbatim}
\end{footnotesize}

%---------------------------->
\subsubsection{Accuracy assessment}

The uncertainty in image classification associated to the distribution of pixels to target land cover classes within the datasets analysed in this work might affect the overall results of mapping. Since this is related to monitoring habitat and vegetation over time, spatial uncertainty should be estimated to avoid biased conclusions. In GRASS GIS, a solution evaluate accuracy is proposed by the chi square test which plots the reject raster map which demonstrates the reject threshold results. The aim is to determine the confidence level at which each cell is categorised correctly and accurately assigned to the correct land cover class. It is the reject threshold map layer, and contains the index to one calculated confidence level for each classified cell in the classified image. The confidence intervals are predefined with the reject map interpreted as values starting from the lowest rejection level (0\% = keep, i.e., the pixels are classified correctly) to the highest (100\% = reject, i.e., low probability of the correctness of the pixel assigned to the target class).Technically, the estimation of accuracy is performed using the scripts presented below. This algorithm applies the chi square test on each discriminant result at various threshold levels of confidence. 

\begin{footnotesize}
\begin{verbatim}
# Mapping rejection probability
d.mon wx1
g.region raster=L_2019_clusters -p
r.colors L_2019_cluster_reject color=soilmoisture -e
d.rast L_2019_cluster_reject
d.grid -g size=00:30:00 color=grey width=0.1 \
fontsize=16 text_color=grey
d.legend raster=L_2019_cluster_reject title="2019" \
title_fontsize=19 \
	font="Helvetica" fontsize=17 bgcolor=white \
	border_color=white
d.out.file output=Italy_2019_reject format=jpg
\end{verbatim}
\end{footnotesize}

%---------------------------->
\subsubsection{RandomForestClassifier}

The RandomForestClassifier (RF Classifier) used for image processing in GRASS GIS belong to the group of Machine Learning (ML) algorithms\hl{ }\citep{598994}. \hl{The fundamental feature of the RandomForestClassifier ML technique is that, when processing the images for training, it creates a complexity of the trees} \citep{990132}. More specifically, \hl{by randomizing the classification, it generates each tree in the ensemble from the training set of categorised pixels using their replacement.}. In GRASS GIS, the following script with defined parameters was has been applied to use RF algorithm for satellite image processing: 

\begin{footnotesize}
\begin{verbatim}
r.learn.train group=L_2019 training_map=training_pixels \
model_name=RandomForestClassifier \
    n_estimators=500 save_model=rf_model.gz --overwrite
# perform prediction using r.learn.predict
r.learn.predict group=L_2019 \
	load_model=rf_model.gz output=rf_classification 
# display
r.colors rf_classification color=bcyr -e
d.mon wx0
d.rast rf_classification
d.grid -g size=00:30:00 color=grey width=0.1 fontsize=16 \
text_color=grey
d.legend raster=rf_classification title="RF 2019" \
title_fontsize=19 font="Helvetica" \
fontsize=17 bgcolor=white border_color=white
d.out.file output=RF_2019 format=jpg --overwrite
\end{verbatim}
\end{footnotesize}

%---------------------------->
\subsubsection{SVC Classifier}

The main approach of the Support Vector Classification SVC \citep{Cortes} for remote sensing data\hl{ }is classifying the new satellite data\hl{ }using the decision process\hl{. }The algorithm of SVC embedded in the GRASS GIS solves this problem through the maximisation of margins of the sample classes within each satellite image by analysing the Digital Numbers (DN) of pixels on the raster scene and inverse regularisation parameter. \hl{In} GRASS GIS, the following scripts. Specifically, the algorithm was used for applying the SVM for image processing and performing prediction using 'r.learn.predict' module and training a SVC model using r.learn.train:

\begin{footnotesize}
\begin{verbatim}
r.learn.train group=L_2023a training_map=training_pixels \
model_name=SVC n_estimators=500 \
save_model=svc_model.gz 
r.learn.predict group=L_2023a load_model=svc_model.gz \
output=svc_classification --overwrite
r.colors svc_classification color=bgyr -e
d.mon wx0
d.rast svc_classification
d.grid -g size=00:30:00 color=grey width=0.1 \
fontsize=16 text_color=grey
d.legend raster=svc_classification title="SVM 2023" \
title_fontsize=19 font="Helvetica" fontsize=17 \
bgcolor=white border_color=white
d.out.file output=SVM_2023 format=jpg --overwrite
\end{verbatim}
\end{footnotesize}

%---------------------------->
\subsubsection{DecisionTree Classifier}

The DecisionTree (DTC) Classifier is\hl{ }based on the supervised learning algorithm. It is considered among one of the most intuitive classifiers\hl{ }and straightforward implementation. The\hl{ }image partition \hl{assigns} the pixels\hl{ }according to the decision rule which evaluates \hl{their} suitability of pixels into a target class\hl{. }For remote sensing approach, the estimation of correctness is computed using the approach of splitting\hl{ }raster dataset into the group categories of pixels\hl{. Each} pixel is defined to be suitable to this specific cluster or not according to the spectral reflectance value. Thus, the algorithm of DecisionTree Classifier computes the probability likelihood of the features representing\hl{ }the \hl{Earth's landscapes}. This considers the value of cells assigned to each class of the classified satellite image as variables:

\begin{footnotesize}
\begin{verbatim}
# train a DTC model using r.learn.train
r.learn.train group=L_2023a training_map=training_pixels \
model_name=DecisionTreeClassifier \
n_estimators=500 save_model=dtc_model.gz
# perform prediction using r.learn.predict
r.learn.predict group=L_2023a load_model=dtc_model.gz \
output=dtc_classification --overwrite
r.colors dtc_classification color=plasma -e
d.mon wx0
d.rast dtc_classification
d.grid -g size=00:30:00 color=grey width=0.1 \
fontsize=16 text_color=grey
d.legend raster=dtc_classification title="DTC 2023" \
	title_fontsize=19 font="Helvetica" fontsize=17 \
	bgcolor=white border_color=white
d.out.file output=DTC_2023 format=jpg --overwrite
\end{verbatim}
\end{footnotesize}

%---------------------------->
\subsubsection{MLPClassifier}

\hl{The} practical solution of GRASS GIS to use the MLPC ANN algorithm \cite{Rosenblatt} consists in using the following script where the name of the model is explicitly defined as MLPClassifier (below an example for the image on 2022, autumn): 

\begin{footnotesize}
\begin{verbatim}
r.learn.train group=L_2022a training_map=training_pixels \
model_name=MLPClassifier n_estimators=500 \
    save_model=mlpc_model.gz --overwrite
# perform prediction using r.learn.predict
r.learn.predict group=L_2022a load_model=mlpc_model.gz \
output=mlpc_classification
# display
r.colors mlpc_classification color=inferno -e
d.mon wx0
d.rast mlpc_classification
d.grid -g size=00:30:00 color=grey width=0.1 fontsize=16 \
text_color=grey
d.legend raster=mlpc_classification title="ANN 2022" \
title_fontsize=19 font="Helvetica" fontsize=17 \
bgcolor=white border_color=white
d.out.file output=MLPC_2022 format=jpg --overwrite
\end{verbatim}
\end{footnotesize}

The results of the MLPClassifer applied to eight Landsat images covering the region of central Apennines including Foreste Casentinesi National Park are presented in figures below.

\section{Results}

The cartographic visualization of the images processed with four methods is presented in the figures above. The land cover types in the northern Italy were identified as following 10 classes according to the existing adopted, modified and generalised classification schemes \citep{DeFioravante}: 1) urban areas, 2) arable lands and agricultural fields, 3) water areas, 4) needle-leaved trees, 5) bare lands and soils; 6) pastures; 7) permanent herbaceous vegetation; 8) croplands; 9) broadleaved trees; 10) shrubland. 

Among the compared outputs regarding the classified images using algorithms of Support Vector Machine (SVM), Multi-layer Perceptron (MLP) Classifier of the Artificial Neural Network (ANN), Decision Tree Classifier (DTC), and Random Forest (RF), the following remarks can be provided. The Random Forest (Figure \ref{fig01}) and SVC Classifiers (Figure \ref{fig02}) differ from the  and Decision Tree (Figure \ref{fig03}) and ANN approaches (Figure \ref{fig04}), due to the different methods in identifying cells within the raster scenes (MLP of ANN and DTC). Hence, the maps generated using SVM method (Figure \ref{fig02}, a) are comparable with those based on the RF Classifier, while MLP Classifier presented a more generalised map (Figure \ref{fig04} using deep learning network approach. 

%\begin{figure}[H]
\begin{figure*}
\centering
\hspace{10pt}\makebox[0.99\linewidth][r]{%
	\begin{subfigure}[b]{.25\textwidth}
		\centering
			\includegraphics[width=0.98\textwidth]{Fig_01a.jpg}
		\caption{20 April 2018}
	\end{subfigure}%
	\begin{subfigure}[b]{.25\textwidth}
		\centering
			\includegraphics[width=0.98\textwidth]{Fig_01b.jpg}
		\caption{22 March 2019}
	\end{subfigure}%
	\begin{subfigure}[b]{.25\textwidth}
		\centering
			\includegraphics[width=0.98\textwidth]{Fig_01c.jpg}
		\caption{14 March 2022}
	\end{subfigure}%
	\begin{subfigure}[b]{.25\textwidth}
		\centering
			\includegraphics[width=0.99\textwidth]{Fig_01d.jpg}
		\caption{17 March 2023}
	\end{subfigure}%
}\\
\vfill \vspace{0.2mm}
\hspace{10pt}\makebox[0.99\linewidth][r]{%
	\begin{subfigure}[b]{.25\textwidth}
		\centering
			\includegraphics[width=0.98\textwidth]{Fig_01e.jpg}
		\caption{27 September 2018}
	\end{subfigure}%
	\begin{subfigure}[b]{.25\textwidth}
		\centering
			\includegraphics[width=0.99\textwidth]{Fig_01f.jpg}
		\caption{14 September 2019}
	\end{subfigure}%
	\begin{subfigure}[b]{.25\textwidth}
		\centering
			\includegraphics[width=0.98\textwidth]{Fig_01g.jpg}
		\caption{6 September 2022}
	\end{subfigure}%
	\begin{subfigure}[b]{.25\textwidth}
		\centering
			\includegraphics[width=0.99\textwidth]{Fig_01h.jpg}
		\caption{3 October 2023}
	\end{subfigure}%
}
\caption{Time series of the multispectral images classified using Random Forest (RF) algorithm.}
\label{fig01}
%\end{figure}
\end{figure*}

Being one of the most efficient inductive learning algorithms of ML in data science in general and remote sensing in particular, SVM has advantages in the classification performances among the ML algorithms due to the effects of dependence distribution of nodes within classes and integrated dependence of all nodes. In this way, it works similar to the Naïve Bayes approach. Hence, SVM is optimal for classification of images where the dependences distribute evenly in classes which is excellently suits to the homogeneous types of landscapes with diversified land cover classes. However, this approach is not optimised for images with high landscape heterogeneity and contrasting landscape types (e.g., sharp gradient between the coastal and mountainous areas).

The advantages of SVM (Figure \ref{fig02} can be explained by its effectivity in high dimensional datasets which is valuable for satellite images that consists thousands of pixels as cells. In our case, 8 multispectral images were classified automatically for each band of Landsat 8-9 OLI/TIRS scenes with 30-m resolution. The target periods cover spring and autumn periods with different vegetation types. Each raster image of the Landsat image contained 7911 rows and 7762 columns of cells, therefore, the automatic approach of processing there data was essential to ensure rapid and accurate data handling. 

Following SVM, the RF classifier also presented optimal results of classification compared to the other approaches in the presented maps. Its superiority is explained by the tuned optimization of the image classification. In the case of RF, using the parameter tuning in GRASS GIS, each node is being split during the construction of a tree with the best split found either from all input pixels of the image or a randomised subset depending on the resolution and extent of the training dataset of the image which depends on the study area and heterogeneity of the visualised landscapes. This randomness aims at decreasing the variance of the RF estimator which allows to overcome overfitting of the individual decision trees with high variance in case of heterogeneous landscapes typical for the Foreste Casentinesi National Park area. The artificially increased randomness of the RF algorithm enables either to decrease the prediction errors of the decision trees or cancel them out, thus making image classification more accurate. 

\begin{figure*}
\centering
\hspace{10pt}\makebox[0.99\linewidth][r]{%
	\begin{subfigure}[b]{.25\textwidth}
		\centering
			\includegraphics[width=0.98\textwidth]{Fig_02a.jpg}
		\caption{20 April 2018}
	\end{subfigure}%
	\begin{subfigure}[b]{.25\textwidth}
		\centering
			\includegraphics[width=0.98\textwidth]{Fig_02b.jpg}
		\caption{22 March 2019}
	\end{subfigure}%
	\begin{subfigure}[b]{.25\textwidth}
		\centering
			\includegraphics[width=0.98\textwidth]{Fig_02c.jpg}
		\caption{14 March 2022}
	\end{subfigure}%
	\begin{subfigure}[b]{.25\textwidth}
		\centering
			\includegraphics[width=0.99\textwidth]{Fig_02d.jpg}
		\caption{17 March 2023}
	\end{subfigure}%
}\\
\vfill \vspace{0.2mm}
\hspace{10pt}\makebox[0.99\linewidth][r]{%
	\begin{subfigure}[b]{.25\textwidth}
		\centering
			\includegraphics[width=0.98\textwidth]{Fig_02e.jpg}
		\caption{27 September 2018}
	\end{subfigure}%
	\begin{subfigure}[b]{.25\textwidth}
		\centering
			\includegraphics[width=0.99\textwidth]{Fig_02f.jpg}
		\caption{14 September 2019}
	\end{subfigure}%
	\begin{subfigure}[b]{.25\textwidth}
		\centering
			\includegraphics[width=0.99\textwidth]{Fig_02g.jpg}
		\caption{6 September 2022}
	\end{subfigure}%
	\begin{subfigure}[b]{.25\textwidth}
		\centering
			\includegraphics[width=0.98\textwidth]{Fig_02h.jpg}
		\caption{3 October 2023}
	\end{subfigure}%
}
\caption{Time series of the multispectral images classified using Support Vector Machine (SVM) algorithm.}\label{fig02}
\end{figure*}

More specifically, the SVM presents high search accuracy compared to other algorithms after a few iterations. Besides, using SVM is supportive for time series analysis where several satellite images are compared since this method is effective when number of dimensions is larger than samples. This is particular valuable for detecting complex land cover patterns in the regions of central Apennines with high heterogeneity of landscapes, such as mixed areas with urban areas and natural spaces. 

\begin{figure*}
\centering
\hspace{10pt}\makebox[0.99\linewidth][r]{%
	\begin{subfigure}[b]{.25\textwidth}
		\centering
			\includegraphics[width=0.98\textwidth]{Fig_03a.jpg}
		\caption{20 April 2018}
	\end{subfigure}%
	\begin{subfigure}[b]{.25\textwidth}
		\centering
			\includegraphics[width=0.98\textwidth]{Fig_03b.jpg}
		\caption{22 March 2019}
	\end{subfigure}%
	\begin{subfigure}[b]{.25\textwidth}
		\centering
			\includegraphics[width=0.99\textwidth]{Fig_03c.jpg}
		\caption{14 March 2022}
	\end{subfigure}%
	\begin{subfigure}[b]{.25\textwidth}
		\centering
			\includegraphics[width=0.99\textwidth]{Fig_03d.jpg}
		\caption{17 March 2023}
	\end{subfigure}%
}\\
\vfill \vspace{0.2mm}
\hspace{10pt}\makebox[0.99\linewidth][r]{%
	\begin{subfigure}[b]{.25\textwidth}
		\centering
			\includegraphics[width=0.98\textwidth]{Fig_03e.jpg}
		\caption{27 September 2018}
	\end{subfigure}%
	\begin{subfigure}[b]{.25\textwidth}
		\centering
			\includegraphics[width=0.98\textwidth]{Fig_03f.jpg}
		\caption{14 September 2019}
	\end{subfigure}%
	\begin{subfigure}[b]{.25\textwidth}
		\centering
			\includegraphics[width=0.98\textwidth]{Fig_03g.jpg}
		\caption{6 September 2022}
	\end{subfigure}%
	\begin{subfigure}[b]{.25\textwidth}
		\centering
			\includegraphics[width=0.99\textwidth]{Fig_03h.jpg}
		\caption{3 October 2023}
	\end{subfigure}%
}
\caption{Time series of the multispectral images classified using DecisionTree algorithm}\label{fig03}
\end{figure*}

\begin{figure*}
\centering
\hspace{10pt}\makebox[0.99\linewidth][r]{%
	\begin{subfigure}[b]{.25\textwidth}
		\centering
			\includegraphics[width=0.98\textwidth]{Fig_04a.jpg}
		\caption{20 April 2018}
	\end{subfigure}%
	\begin{subfigure}[b]{.25\textwidth}
		\centering
			\includegraphics[width=0.98\textwidth]{Fig_04b.jpg}
		\caption{22 March 2019}
	\end{subfigure}%
	\begin{subfigure}[b]{.25\textwidth}
		\centering
			\includegraphics[width=0.99\textwidth]{Fig_04c.jpg}
		\caption{14 March 2022}
	\end{subfigure}%
	\begin{subfigure}[b]{.25\textwidth}
		\centering
			\includegraphics[width=0.99\textwidth]{Fig_04d.jpg}
		\caption{17 March 2023}
	\end{subfigure}%
}\\
\vfill \vspace{0.2mm}
\hspace{10pt}\makebox[0.99\linewidth][r]{%
	\begin{subfigure}[b]{.25\textwidth}
		\centering
			\includegraphics[width=0.98\textwidth]{Fig_04e.jpg}
		\caption{27 September 2018}
	\end{subfigure}%
	\begin{subfigure}[b]{.25\textwidth}
		\centering
			\includegraphics[width=0.98\textwidth]{Fig_04f.jpg}
		\caption{14 September 2019}
	\end{subfigure}%
	\begin{subfigure}[b]{.25\textwidth}
		\centering
			\includegraphics[width=0.98\textwidth]{Fig_04g.jpg}
		\caption{6 September 2022}
	\end{subfigure}%
	\begin{subfigure}[b]{.25\textwidth}
		\centering
			\includegraphics[width=0.98\textwidth]{Fig_04h.jpg}
		\caption{3 October 2023}
	\end{subfigure}%
}
\caption{Time series of the multispectral ANN-based MLPC algorithm}\label{fig04}
\end{figure*}

\section{Conclusions}

This work demonstrated the examples of the supervised training in ML approaches to geospatial applications. The research evaluated ML algorithms of image classification for improved workflow of image processing and tested their applicability in cartographic works. Four algorithms of ML supervised training were compared using GRASS GIS modules that are based Python's library Scikit-Learn. Based on the maps generated from the processed Landsat satellite images, four distinct ML approaches\hl{ }were compared both for mathematical foundation and graphical output. 

Geospatial datasets are important information source for analysis of landscape dynamics. Using spaceborne data enables to model biodiversity changes in a long-term scale. Such time series analysis enables to compare current and previous state of the landscapes using cartographic data analysis and visualization which can be used for environmental monitoring, assessment of forest health and prognosis of possible further development in the near future. Landscape dynamics includes a complexity of both human-induced and climate-related processes including land abandonment as its major side effect. Landscape dynamics is triggered by a consequence of the expansion of forest area and the habitat homogenization. Early detection of land cover changes as parts of landscape needs a timely recognition of subtle changes in the forests and involves fieldwork monitoring. One of the important approaches for identification landscape dynamics in land cover types is the analysis of changes and the rate of exposure of landscapes to climate and environmental effects using multi-temporal geospatial data. 

The consequences of various ML parameters on the cartographic outputs are analysed in this work, such as speed and accuracy, randomness of nodes, analytical determination of the output weights, and dependence distribution of pixels for each algorithm. Supervised learning models of GRASS GIS were tested and compared. The results shown that the most time-consuming algorithms was noted as SVC classification, while the fastest results were achieved by the Decision Tree approach to image processing and the best results are achieved by RF Classifier. Though each algorithms was developed to serve different objectives of ML applications in RS data processing, their technical implementation and practical purposes present valuable approaches to cartographic data processing and image analysis. 

To conclude, the ML applications can be effectively used for environmental monitoring and forest mapping. Using ML approach, the automation of geospatial data processing increases according to the choice of diverse ML instruments. The particular case of this work presented an example of the ML modules of GRASS GIS to process the multispectral Landsat images. Several algorithms of ML were tested and compared for visualization of landscapes using image partition into land cover classes automatically. The effectiveness of the remote sensing (RS) data for environmental forest monitoring is explained by the optical properties of satellite images: the features of land cover types on the Earth have different spectral reflectance values. This enables to identify them using digital signatures of the satellite images. Time series analysis of such scenes enables to detect changes and recognise modified patches within the landscape structure. Therefore, extraction of these features from the RS signals recorded during the Earth monitoring from space missions (such as Landsat) provides an essential tool to detect landscape dynamics. 


\bibliographystyle{actapoly}
%\bibliographystyle{apalike}
\bibliography{References_ML.bib}

\end{document}
